\subsection{共轭定理}

定理 2. —— 设 B 是一个半单环，Z 是它的中心；假设 B 是一个有限生成的 Z-模。设 A 是一个右阿廷环，f 与 g 是从 B 到 A 的环同态；用 \(f_{Z}\) 与 \(g_{Z}\) 表示 f 与 g 在 Z 上的限制。下列性质等价：

\begin{enumerate}
\def\labelenumi{(\roman{enumi})}
\item
  存在 A 的一个内自同构 \(\theta\)，使得 \(g = \theta \circ f\)。
\item
  存在 A 的一个内自同构 \(\theta\)，使得 \(g_{\mathbb{Z}} = \theta \circ f_{\mathbb{Z}}\)。
\end{enumerate}

由于环 A 是右阿廷的，\(A_{d}\) 是一个有限长度的右 A-模（VIII, 第 6 页，定理 1）。因此 \(A^{f}\) 与 \(A^{g}\) 是有限长度的 (B, A)-双模。由引理 1（VIII, 第 254 页），断言 (i) 表示 \(A^{f}\) 与 \(A^{g}\) 是同构的 (B, A)-双模，而断言 (ii) 表示它们是同构的 (Z, A)-双模。于是 (i) 与 (ii) 的等价性由引理 2（VIII, 第 254 页）得出。

推论. —— 设 A 与 B 是域 K 上的代数。假设 B 是中心的、单的且为有限次数的，而 A 是右阿廷的。设 f 与 g 是从 B 到 A 的 K-代数同态。存在 A 的一个内自同构 \(\theta\)，使得 \(g = \theta \circ f\)。

在定理 2 的记号下，我们有 \(\mathrm{Z} = \mathrm{K}\)，因此 \(f_{\mathrm{Z}} = 1_{\mathrm{Z}} = g_{\mathrm{Z}}\)。

定理 3（斯科伦–诺特定理）. —— 设 A 与 B 是单 K-代数，Z(A) 与 Z(B) 是它们的中心。假设代数 B 在 K 上是有限次数的，且代数 \(Z(A) \otimes_{K} Z(B)\) 是一个域（特别地，当 A 或 B 是中心的时候即是这种情形）。设 f 与 g 是从 B 到 A 的 K-代数同态。存在 A 的一个内自同构 \(\theta\)，使得 \(g = \theta \circ f\)。

由 VIII, 第 254 页的引理 1，只需证明 (B, A)-双模 \(A^{f}\) 与 \(A^{g}\) 同构。现在，我们可以把 \(A^{f}\) 与 \(A^{g}\) 视为代数 \(C = B \otimes_{K} A^{o}\) 上的左模，而由 VIII, 第 221 页的命题 7，该代数是单的。作为右 A-模，\(A^{f}\) 与 \(A^{g}\) 同构于 \(A_{d}\)，因而由于环 A 是单的而具有有限长度（VIII, 第 121 页，推论 1）。更有甚者，\(A^{f}\) 与 \(A^{g}\) 是有限长度的 C-模。设 S 是一个单 C-模；存在严格正的整数 m 与 n，使得 \(A^{f}\) 同构于 \(S^{m}\)，而 \(A^{g}\) 同构于 \(S^{n}\)。于是右 A-模 S 具有非零的有限长度。由于 \(A^{f}\) 与 \(A^{g}\) 的基础右 A-模是同构的，它们具有相同的长度；因此我们有 m = n，从而 C-模 \(A^{f}\) 与 \(A^{g}\) 同构。

推论 1. —— 设 A 是 K 上的一个中心单代数，L 是 K 的一个有限次数的扩张。若 f 与 g 是从 L 到 A 的 K-代数同态，则存在 A 的一个内自同构 \(\theta\)，使得 \(g = \theta \circ f\)。

推论 2. —— 设 A 是 K 上的一个中心单代数，L 是 A 的一个子代数且是一个体。从 L 到 A 的每个 K-代数同态都可以扩张为 A 的一个内自同构。

推论 3. —— 设 D 是 K 上的一个有限次数的体，其中心为 K。D 的每个元素在 K 上都是代数的。设 x 与 y 是 D 的元素；存在 \(D^{*}\) 的一个元素 a 使得 \(y = axa^{-1}\)，当且仅当 x 与 y 在 K 上有相同的最小多项式。

第一个断言由 V, §3, No.~1, 第 17 页的推论 1 得出。

假设存在一个元素 \(a\)（\(\mathrm{D}^*\) 的），使得 \(y = axa^{-1}\)；对每个多项式 \(\mathrm{P}\)（\(\mathrm{K}[\mathrm{X}]\) 的），我们有 \(\mathrm{P}(y) = a\mathrm{P}(x)a^{-1}\)，特别地，\(\mathrm{P}(x) = 0\) 当且仅当 \(\mathrm{P}(y) = 0\)。因此，\(x\) 与 \(y\) 在 \(\mathrm{K}\) 上有相同的最小多项式（V, §3, No.~1, 第 16 页，定理 1）。

反之，假设 x 与 y 有相同的最小多项式。由同上，存在从 K{[}x{]} 到 K{[}y{]} 的一个 K-同构 u，使得 \(u(x) = y\)，且 K{[}x{]} 是一个域。由推论 2，u 扩张为 D 的一个内自同构 \(\theta : z \mapsto aza^{-1}\)，因此我们有 \(y = \theta(x) = axa^{-1}\)。

命题 1. —— 设 A 是 K 上的一个有限次数的中心单代数。设 B 是一个 K-代数，f 与 g 是从 B 到 A 的代数同态。下列性质等价：

\begin{enumerate}
\def\labelenumi{(\roman{enumi})}
\item
  存在 A 的一个内自同构 \(\theta\)，使得 \(g = \theta \circ f\)。
\item
  作为左 B-模，\(\mathrm{A}^f\) 与 \(\mathrm{A}^g\) 同构。
\end{enumerate}

由引理 1（VIII, 第 254 页），性质 (i) 等价于 \(A^{f}\) 与 \(A^{g}\) 作为 (B, A)-双模同构这一事实。由于 A 在 K 上是有限维的，\(A^{f}\) 与 \(A^{g}\) 是有限长度的 B-模。由于 A 的中心等于 K，(i) 与 (ii) 的等价性由应用于 \((\mathrm{A}^{\circ}, \mathrm{B}^{\circ})\)-双模 \(A^{f}\) 与 \(A^{g}\) 的 VIII, 第 254 页的引理 2 得出。

